\documentclass[12pt,a4paper]{article}

\usepackage[utf8]{inputenc}
\usepackage[T1]{fontenc}
\usepackage[french]{babel}
\usepackage{geometry}
\geometry{margin=2.5cm}
\usepackage{graphicx}
\usepackage{amsmath}
\usepackage{amssymb}
\usepackage{booktabs}
\usepackage{array}
\usepackage{xcolor}
\usepackage{listings}
\usepackage{hyperref}
\usepackage{enumitem}
\usepackage{float}
\usepackage{caption}
\usepackage{subcaption}
\usepackage{tikz}
\usetikzlibrary{shapes.geometric, arrows.meta, positioning, fit, backgrounds, calc}
\usepackage{pgfplots}
\pgfplotsset{compat=1.18}
\usepackage{eurosym}
\usepackage{microtype}
\usepackage{lmodern}
\usepackage{fancyhdr}
\usepackage{titlesec}
\usepackage{parskip}

\setlength{\parskip}{6pt}
\setlength{\parindent}{0pt}

\pagestyle{fancy}
\fancyhf{}
\fancyhead[L]{\small Agent de Preprocessing Data Science}
\fancyhead[R]{\small \thepage}
\renewcommand{\headrulewidth}{0.4pt}

\definecolor{codeblue}{RGB}{0,90,160}
\definecolor{codegray}{RGB}{240,240,240}
\definecolor{codered}{RGB}{180,30,30}

\lstset{
  backgroundcolor=\color{codegray},
  basicstyle=\small\ttfamily,
  keywordstyle=\color{codeblue}\bfseries,
  commentstyle=\color{gray}\itshape,
  stringstyle=\color{codered},
  breaklines=true,
  frame=single,
  rulecolor=\color{gray!50},
  numbers=left,
  numberstyle=\tiny\color{gray},
  numbersep=5pt,
  tabsize=2,
  captionpos=b
}

\hypersetup{
  colorlinks=true,
  linkcolor=codeblue,
  urlcolor=codeblue,
  citecolor=codeblue
}

\tikzset{
  agent/.style={
    rectangle, rounded corners=4pt, draw=codeblue, fill=codeblue!10,
    text width=4.2cm, minimum height=1.1cm, align=center, font=\small\bfseries
  },
  node_box/.style={
    rectangle, rounded corners=3pt, draw=gray!70, fill=white,
    text width=3.6cm, minimum height=0.8cm, align=center, font=\footnotesize
  },
  interrupt/.style={
    diamond, draw=codered, fill=codered!10,
    aspect=2.5, minimum height=0.8cm, align=center, font=\footnotesize\bfseries
  },
  arrow/.style={-{Stealth[length=6pt]}, thick, draw=gray!70},
  redarrow/.style={-{Stealth[length=6pt]}, thick, draw=codered},
}

\title{
  \vspace{-1cm}
  \Large\textbf{Agent de Preprocessing Data Science}\\[0.4cm]
  \large Conception d'un pipeline de préparation de données semi-automatisé\\[0.6cm]
  \normalsize Rapport technique
}
\author{Ibrahim Youssouf \textsc{abdelatif}}
\date{Mars 2026}

\begin{document}

\maketitle
\thispagestyle{empty}

\vspace{0.5cm}

\begin{abstract}
Ce rapport documente la conception d'un pipeline de preprocessing automatisé piloté
par GPT-5-mini. L'outil prend en entrée un fichier CSV avec une colonne cible et
produit des jeux train/test prêts pour la modélisation, accompagnés d'un rapport de
qualité décomposé. L'architecture repose sur LangGraph pour orchestrer quatre agents
spécialisés et trois points d'interruption humaine. Le fil conducteur de l'ensemble
des choix d'implémentation est une séparation stricte entre ce qui peut être codé
de façon déterministe et ce qui nécessite génuinement du raisonnement contextuel.
\end{abstract}
\newpage
\vspace{0.5cm}
\tableofcontents
\newpage

% ============================================================
\section{Motivation}
% ============================================================

Le preprocessing consomme en pratique la majorité du temps d'un projet ML, et c'est
là que les erreurs les plus tenaces s'introduisent --- non par négligence, mais par
accumulation de micro-décisions répétitives. La fuite de données du test vers le train,
l'imputation calculée sur l'ensemble complet, l'encodage inconsistant entre les splits :
ces bugs sont aussi courants chez les praticiens expérimentés que chez les débutants,
précisément parce qu'ils ne résultes pas d'une incompréhension mais d'une inattention.

L'objectif du projet est d'automatiser ces décisions sans pour autant construire une
boîte noire. Deux contraintes ont guidé l'implémentation depuis le début.

La première : tout ce qui est non ambigu doit être codé en dur. Supprimer une colonne
à 100\% de valeurs manquantes ne dépend pas du domaine. Ces règles sont déterministes,
testables unitairement, et ne nécessitent pas de modèle de langage.

La seconde : le modèle de langage n'intervient que là où une règle générique serait
arbitraire. Choisir entre \texttt{one\_hot\_encoding} et \texttt{frequency\_encoding}
pour une colonne à cardinalité moyenne n'a pas de réponse universelle --- cela dépend
du domaine, de la sémantique des catégories, et du type de modèle cible. C'est
précisément le genre de raisonnement où un LLM nourri du contexte du dataset apporte
quelque chose qu'une heuristique codée en dur ne peut pas offrir.

% ============================================================
\section{Architecture générale}
% ============================================================

\subsection{Choix du framework : LangGraph}

Plusieurs options ont été évaluées. LangChain en mode chain classique couvre les flux
linéaires, mais ne gère pas nativement les boucles de rejet humain ni la persistance
d'état entre deux interruptions. Prefect et Airflow sont de bons outils d'orchestration
mais n'ont pas de primitives pour le human-in-the-loop. Du Python pur aurait été la
solution la plus simple, sauf que la gestion manuelle des checkpoints et des reprises
après interruption devient fragile dès que le pipeline dépasse trois ou quatre étapes.

LangGraph s'impose par deux caractéristiques concrètes. Le \texttt{StateGraph} partage
un dictionnaire d'état mutable entre tous les nœuds, ce qui évite de faire transiter
des données par des variables globales ou des fichiers intermédiaires. Le \texttt{interrupt()}
suspend l'exécution en sérialisant l'état complet dans un checkpointer SQLite : si
l'utilisateur ferme son terminal entre deux validations, la reprise se fait en
réinstanciant le pipeline avec le même \texttt{thread\_id} --- l'état est restauré
exactement depuis le dernier nœud exécuté.

\subsection{Vue d'ensemble du graphe}

Le pipeline est composé de onze nœuds et trois points d'interruption humaine
(figure~\ref{fig:pipeline}). Chaque interruption soumet les propositions de l'agent
courant à l'utilisateur ; un rejet renvoie au nœud précédent sans compteur de retries,
laissant intentionnellement l'humain converger à son rythme.

\begin{figure}[H]
\centering
\begin{tikzpicture}[node distance=0.55cm and 0.3cm]

  \node[node_box, fill=gray!10] (start) {START};
  \node[node_box, below=of start] (analyze) {analyze};
  \node[node_box, below=of analyze] (split) {split\_data};
  \node[interrupt, below=of split] (hr_domain) {Revue\\domaine};
  \node[node_box, below=of hr_domain] (propose_t) {propose\_transformations};
  \node[interrupt, below=of propose_t] (hr_trans) {Revue\\transforms};
  \node[interrupt, right=1.8cm of hr_trans, fill=orange!10, draw=orange] (hr_leak) {Alerte\\leakage};
  \node[node_box, below=of hr_trans] (exec_t) {execute\_transformations};
  \node[node_box, below=of exec_t] (propose_o) {propose\_outlier\_treatment};
  \node[interrupt, below=of propose_o] (hr_out) {Revue\\outliers};
  \node[node_box, below=of hr_out] (exec_o) {execute\_outlier\_treatment};
  \node[node_box, below=of exec_o] (scaling) {apply\_scaling};
  \node[node_box, below=of scaling] (report) {generate\_report};
  \node[node_box, below=of report, fill=gray!10] (end) {END};

  \draw[arrow] (start) -- (analyze);
  \draw[arrow] (analyze) -- (split);
  \draw[arrow] (split) -- (hr_domain);
  \draw[arrow] (hr_domain) -- node[right, font=\tiny]{approve} (propose_t);
  \draw[arrow] (propose_t) -- (hr_trans);
  \draw[arrow] (hr_trans) -- node[right, font=\tiny]{approve} (exec_t);
  \draw[arrow] (exec_t) -- (propose_o);
  \draw[arrow] (propose_o) -- (hr_out);
  \draw[arrow] (hr_out) -- node[right, font=\tiny]{approve} (exec_o);
  \draw[arrow] (exec_o) -- (scaling);
  \draw[arrow] (scaling) -- (report);
  \draw[arrow] (report) -- (end);

  \draw[redarrow] (hr_domain.west) -- ++(-1.2,0) |- node[left, font=\tiny, pos=0.25]{reject} (analyze.west);
  \draw[redarrow] (hr_trans.west) -- ++(-1.4,0) |- node[left, font=\tiny, pos=0.25]{reject} (propose_t.west);
  \draw[redarrow] (hr_out.west) -- ++(-1.2,0) |- node[left, font=\tiny, pos=0.25]{reject} (propose_o.west);

  \draw[arrow] (hr_trans.east) -- node[above, font=\tiny]{leakage} (hr_leak.west);
  \draw[arrow] (hr_leak.south) |- (exec_t.east);

\end{tikzpicture}
\caption{Graphe d'exécution LangGraph. Les losanges rouges sont des points d'interruption
humaine. Les flèches rouges représentent les boucles de rejet.}
\label{fig:pipeline}
\end{figure}

\subsection{Le \texttt{PipelineState} comme contrat entre agents}

Tous les nœuds communiquent exclusivement via un \texttt{TypedDict} appelé
\texttt{PipelineState}. Ce choix architectural a une conséquence directe sur la
testabilité : chaque agent devient une fonction pure dont les sorties ne dépendent
que de l'état en entrée, sans effets de bord cachés ni dépendance à des variables
externes. Les champs principaux sont les chemins des fichiers CSV à chaque étape
(\texttt{df\_path}, \texttt{train\_path}, \texttt{test\_path}), les statistiques
du dataset initial (\texttt{df\_info}), le contexte métier inféré (\texttt{domain\_context}),
les propositions en attente de validation (\texttt{proposals}), et les artefacts de
preprocessing ajustés sur le train (\texttt{fitted\_artifacts}).

Le champ \texttt{steps\_done} utilise l'annotation \texttt{Annotated[list, operator.add]}
de LangGraph, qui accumule automatiquement les contributions de chaque nœud sans
écrasement. Cela produit un historique complet et ordonné de toutes les décisions
prises au cours du pipeline.

% ============================================================
\section{Agent 1 : Analyse et split (BaseAgent)}
% ============================================================

\subsection{Audit des types}

Pandas ne garantit pas une inférence correcte des types au chargement d'un CSV. Une
colonne numérique stockée comme chaîne de caractères passera inaperçue et sera traitée
comme variable catégorique dans toute la suite du pipeline --- avec des conséquences
silencieuses sur l'encodage, l'imputation, et le scaling.

La détection repose sur un essai vectorisé : pour chaque colonne de type \texttt{object},
\texttt{pd.to\_numeric(..., errors="coerce")} est appliqué et le ratio de valeurs
convertibles est calculé. Un seuil à 95\% signale la colonne comme mal typée.
Le même principe s'applique aux dates avec \texttt{pd.to\_datetime}. Les colonnes
entières ne contenant que $\{0, 1\}$ sont marquées booléennes ; celles avec $\leq 10$
valeurs uniques et un écart inférieur à deux fois leur cardinalité sont marquées
ordinales.

\subsection{Détection des colonnes triviales}

Cinq catégories sont éliminables sans ambiguïté, indépendamment du domaine :

\begin{table}[H]
\centering
\begin{tabular}{lll}
\toprule
\textbf{Catégorie} & \textbf{Critère} & \textbf{Action} \\
\midrule
100\% NaN & toutes valeurs manquantes & suppression \\
Constante & une seule valeur unique & suppression \\
Quasi-constante & valeur dominante $\geq 99{,}5\%$ & suppression \\
Identifiant unique & cardinalité = nb lignes, type objet & suppression \\
Doublon exact & égale bit-à-bit à une autre colonne & suppression \\
\bottomrule
\end{tabular}
\caption{Règles de suppression des colonnes triviales.}
\end{table}

\subsection{Inférence du domaine (LLM)}

Le LLM reçoit les noms de colonnes, un échantillon de cinq lignes, les statistiques
par colonne, et le type de tâche. Il produit un objet JSON valide contre le schéma
Pydantic \texttt{DomainAndAnomaliesResponse}, qui contient le domaine métier inféré,
une description en une ou deux phrases, les contraintes métier (par exemple \textit{age >= 0}),
les colonnes potentiellement sensibles, et les anomalies sémantiques.

Ce dernier point mérite une explication. La détection d'anomalies sémantiques ne peut
pas être codifiée par des règles génériques : une valeur $-1$ dans une colonne
\textit{age} est une sentinelle évidente, mais $-1$ dans une colonne de coordonnée
GPS est parfaitement valide. Une règle qui détecterait l'une manquerait l'autre.
Un second appel produit un diagnostic structuré du dataset (problèmes identifiés,
recommandations), qui sert de contexte aux agents de transformation et de rapport.

\subsection{Split train/test : pourquoi avant les transformations}

Le split est effectué immédiatement après l'analyse, avant toute transformation.
La raison est simple : ajuster un scaler ou calculer une médiane d'imputation sur
l'ensemble complet puis les appliquer aux deux splits constitue une fuite de données.
En splitant en premier, on garantit que tout paramètre de preprocessing est estimé
exclusivement sur le train.

Une nuance mérite d'être explicitée. La phase d'analyse lit le dataset complet pour
produire le contexte domaine et le diagnostic LLM --- ce qui est acceptable, car ces
informations sont de nature exploratoire. En revanche, les statistiques qui pilotent
les décisions de transformation (skewness, ratio d'outliers, quasi-constance) sont
systématiquement recalculées sur le seul jeu train dans \texttt{propose\_transformations},
après le split, de sorte qu'aucune statistique issue du test n'influence les paramètres
de preprocessing.

Quatre stratégies de split sont disponibles : \textbf{auto} (stratifié pour la
classification, aléatoire pour la régression), \textbf{stratified}
(\texttt{StratifiedShuffleSplit} sur la cible), \textbf{temporal} (les derniers
$N\%$ du dataset trié par une colonne temporelle forment le test), et \textbf{group}
(\texttt{GroupShuffleSplit} sur une colonne de groupe, garantissant qu'aucun individu
n'apparaît simultanément dans les deux splits).

% ============================================================
\section{Agent 2 : Transformations (TransformationAgent)}
% ============================================================

\subsection{Deux couches de décision}

L'agent de transformation distingue deux catégories de colonnes. Pour celles dont le
traitement est non ambigu, des règles déterministes s'appliquent directement. Pour les
colonnes ambiguës --- typiquement celles à cardinalité intermédiaire ou dont le domaine
métier conditionne le choix d'encodage --- le LLM est sollicité, et sa proposition est
soumise à validation humaine avant exécution.

Cette organisation a une conséquence pratique sur la traçabilité : chaque décision est
taggée avec sa source (\texttt{rule} ou \texttt{llm}), son niveau de confiance
(\texttt{high}, \texttt{medium}, \texttt{low}), et une justification textuelle.
Le rapport final inclut un \textit{confidence map} qui récapitule l'ensemble de ces
métadonnées de décision.

\begin{figure}[H]
\centering
\begin{tikzpicture}[node distance=0.6cm and 0.5cm,
  decision/.style={diamond, draw=codeblue, fill=codeblue!8, aspect=2.5,
    minimum height=0.7cm, align=center, font=\scriptsize},
  action/.style={rectangle, rounded corners=2pt, draw=gray!70, fill=white,
    text width=2.8cm, minimum height=0.65cm, align=center, font=\scriptsize},
  arr/.style={-{Stealth[length=5pt]}, thick, draw=gray!60}
]
  \node[decision] (d1) {NaN $>$ 50\%?};
  \node[action, right=1.2cm of d1] (drop) {drop\_column [rule]};
  \node[decision, below=0.8cm of d1] (d2) {Mal typé?};
  \node[action, right=1.2cm of d2] (cast) {cast / extract\_datetime [rule]};
  \node[decision, below=0.8cm of d2] (d3) {Cardinalité?};
  \node[action, right=1.2cm of d3] (ohe) {one\_hot $\leq$10 [rule]};
  \node[action, below=0.5cm of ohe] (freq) {freq. enc. $>$50 [rule]};
  \node[action, below=0.5cm of freq] (llm) {LLM 10--50 [llm]};
  \draw[arr] (d1) -- node[above, font=\tiny]{oui} (drop);
  \draw[arr] (d1) -- node[left, font=\tiny]{non} (d2);
  \draw[arr] (d2) -- node[above, font=\tiny]{oui} (cast);
  \draw[arr] (d2) -- node[left, font=\tiny]{non} (d3);
  \draw[arr] (d3) -- node[above, font=\tiny]{$\leq$10} (ohe);
  \draw[arr] (d3.south) -- ++(0,-0.4) -| node{} (freq.west);
  \draw[arr] (d3.south) -- ++(0,-0.4) -| node{} (llm.west);
\end{tikzpicture}
\caption{Arbre de décision par colonne dans le TransformationAgent.}
\label{fig:transform_decision}
\end{figure}

\subsection{Ordre des opérations}

L'application des transformations suit un ordre contraint :

\begin{enumerate}
  \item Corrections de types (\texttt{cast\_numeric}, \texttt{extract\_datetime}).
  \item Imputation (\texttt{impute\_mean/median/mode}) et nettoyage de lignes/colonnes.
  \item Encodage catégorique (\texttt{one\_hot}, \texttt{ordinal}, \texttt{frequency}).
  \item Scaling --- délibérément reporté après le traitement des outliers.
\end{enumerate}

Le report du scaling est une décision d'ordre : standardiser avant le clipping biaiserait
la moyenne et l'écart-type calculés par le scaler, puisque les outliers influeraient sur
ces paramètres. L'ordre correct est clipping d'abord, scaling ensuite.

\subsection{Garanties train/test pour chaque transformation}

Chaque artefact de preprocessing --- imputation, encodeur, scaler, bornes d'outliers ---
est ajusté sur le train uniquement puis appliqué au test avec les paramètres fixes.
Le tableau~\ref{tab:leakage} récapitule ces garanties par type de transformation.

\begin{table}[H]
\centering
\begin{tabular}{lll}
\toprule
\textbf{Transformation} & \textbf{Ajusté sur} & \textbf{Appliqué sur} \\
\midrule
Imputation médiane / mode / moyenne & train & train + test \\
One-hot encoding (catégories sklearn) & train & train + test (catégories inconnues = zéros) \\
Ordinal encoder & train & train + test \\
Frequency encoder & train & train + test \\
StandardScaler / RobustScaler & train & train + test \\
Log / sqrt transform & train & train + test \\
Bornes IQR (outliers) & train & train + test \\
\bottomrule
\end{tabular}
\caption{Étanchéité train/test : aucun paramètre n'est estimé sur le test.}
\label{tab:leakage}
\end{table}

\subsection{Détection de corrélation élevée avec la cible}

Après les propositions de transformation, l'agent calcule la corrélation de Pearson
entre chaque feature numérique et la cible. Toute corrélation supérieure à $0{,}95$
déclenche une alerte et une interruption humaine supplémentaire. Ce seuil détecte les
cas où une feature encode directement la variable cible (identifiant de transaction
corrélé parfaitement au montant), mais peut aussi signaler une redondance légitime
(\texttt{petal\_length} et \texttt{petal\_width} dans Iris). L'humain reste seul juge
de la nature du signal.

\subsection{Multicolinéarité inter-features}

L'agent calcule la matrice de corrélation de Pearson pour les features numériques et
le V de Cramer pour les paires catégoriques. Les paires au-delà de $0{,}75$ sont
signalées en sévérité \texttt{high} ; au-delà de $0{,}90$, en \texttt{very\_high}.
Pour chaque paire, des actions suggérées sont produites : suppression de la feature
la moins informative, ACP, VIF itératif, ou conservation pour les modèles basés
sur les arbres, qui tolèrent la colinéarité mieux que les modèles linéaires.

\subsection{Validation des propositions LLM}

Deux vérifications s'appliquent systématiquement avant de soumettre les propositions à
l'utilisateur. Premièrement, toute proposition \texttt{drop\_column} émise par le LLM
est rejetée : seules les règles déterministes peuvent supprimer des colonnes, ce qui
évite une perte de données non traçable. Deuxièmement, toute proposition référençant
une colonne absente du DataFrame est rejetée silencieusement avec un log d'avertissement.

% ============================================================
\section{Agent 3 : Outliers (OutlierAgent)}
% ============================================================

\subsection{Choix de la méthode de détection}

Le choix entre IQR et Z-score n'est pas trivial. Le Z-score suppose une distribution
approximativement normale : appliqué à une distribution exponentielle ou log-normale,
il détecte un nombre d'outliers largement supérieur au réel. L'agent suit le protocole
suivant pour chaque colonne :

\begin{enumerate}
  \item \textbf{Coefficient de bimodalité de Sarle :} $BC = (\text{skewness}^2 + 1) / (\text{kurtosis} + 3)$.
  Si $BC > 0{,}555$, la distribution est bimodale : ni IQR ni Z-score n'est fiable.
  La colonne est conservée telle quelle.
  \item \textbf{Test de D'Agostino-Pearson} (pour $n \geq 20$) : si $p > 0{,}05$, la
  distribution est approximativement normale, le Z-score est utilisé ; sinon, IQR.
  \item \textbf{Heuristique sur le skewness} (pour $n < 20$) : $|\text{skewness}| < 1{,}0$
  implique Z-score, sinon IQR.
\end{enumerate}

Le test de D'Agostino-Pearson est préféré au test de Shapiro-Wilk, qui perd en puissance
statistique pour $n > 5000$ --- une taille de dataset courante en preprocessing réel.

\subsection{Logique de traitement déterministe}

Le traitement est déterminé par règles avant tout recours au LLM :

\begin{table}[H]
\centering
\begin{tabular}{p{4.5cm}ll}
\toprule
\textbf{Contexte} & \textbf{Action} & \textbf{Confiance} \\
\midrule
Colonne marquée comme sémantique & replace\_with\_nan & high \\
Dataset déséquilibré + peu d'outliers & clip & high \\
Dataset déséquilibré + beaucoup d'outliers & keep & medium \\
$< 1\%$ outliers, classification & remove & high \\
$< 1\%$ outliers, régression & clip & high \\
$1\%$ à $5\%$ outliers & clip & medium \\
$> 5\%$ outliers & keep & low \\
Distribution bimodale & keep & high \\
\bottomrule
\end{tabular}
\caption{Règles déterministes de traitement des outliers selon le contexte.}
\end{table}

Un cas particulier mérite d'être souligné. Supprimer des outliers dans un dataset où
la classe minoritaire représente moins de 15\% des observations peut l'éliminer
entièrement. Lorsque le ratio majoritaire dépasse $0{,}85$, toute suppression est
automatiquement remplacée par du clipping, indépendamment des autres critères.

\subsection{Rôle du LLM}

Le LLM n'intervient que pour les colonnes dont la confiance est \texttt{medium},
généralement celles dont le pourcentage d'outliers est compris entre 1\% et 5\% et
dont le domaine métier peut justifier de les conserver. Un salaire de 500 000~\euro{}
est un outlier statistique sur une population générale, mais pas dans une base de
données de dirigeants de grandes entreprises. C'est ce genre de distinction ---
impossible à coder en dur --- que le LLM apporte ici.

% ============================================================
\section{Agent 4 : Rapport (ReportAgent)}
% ============================================================

\subsection{Score de qualité décomposé}

Le rapport inclut un score sur 100 points calculé entièrement de façon déterministe à
partir de l'état final du pipeline :

\begin{align}
S_{\text{total}} = S_{\text{complétude}} + S_{\text{suppressions}} + S_{\text{types}}
  + S_{\text{leakage}} + S_{\text{outliers}}
\end{align}

Le LLM n'intervient pas dans ce calcul --- il génère uniquement le texte narratif du
rapport. Cela garantit que le score est reproductible : deux exécutions identiques
produisent exactement le même score.

\subsection{Évaluation baseline}

Pour quantifier l'apport du preprocessing, l'agent entraîne deux modèles sur les données
finales. Un \texttt{DummyClassifier} (stratégie \texttt{most\_frequent}) ou
\texttt{DummyRegressor} (moyenne) sert de baseline triviale. Une \texttt{LogisticRegression}
ou une \texttt{Ridge} constitue le modèle simple. Le delta entre les deux sur le jeu
test mesure si les données préprocessées sont exploitables par un modèle linéaire : un
delta négatif signale invariablement un problème de preprocessing résiduel (colonnes
non numériques, NaN non imputés).

\subsection{Vérification des contraintes métier}

Les contraintes inférées par le LLM à l'étape 1 (par exemple \textit{age >= 0}) sont
vérifiées sur les données finales par parsing de regex. Chaque violation est rapportée
avec le nombre et le pourcentage de lignes concernées, ce qui ferme la boucle entre
le contexte métier inféré au début du pipeline et la qualité effective du résultat.

% ============================================================
\section{Infrastructure technique}
% ============================================================

\subsection{Cache LLM et reproductibilité}

Le client \texttt{GPTClient} intègre un cache SQLite persistant dont la clé est un
hash SHA-256 du triplet (modèle, paramètres, messages). Un appel identique --- même
dataset, même version de prompt --- ne déclenche aucun appel API lors d'une ré-exécution.
Cela réduit les coûts de développement de façon significative, et garantit la
reproductibilité des tests : les résultats ne varient pas entre deux exécutions sur le
même dataset, tant que les prompts restent inchangés.

\subsection{Résilience face aux sorties malformées}

Chaque réponse LLM est validée contre un schéma Pydantic strict. En cas d'échec ---
JSON tronqué, champ manquant, action non répertoriée parmi les 17 autorisées ---
le système tente de réparer le JSON (fermeture des crochets et accolades manquants),
puis renvoie le message d'erreur de validation au modèle pour correction automatique.
Ce cycle se répète jusqu'à trois fois avec un backoff exponentiel ($2^{\text{essai}}$
secondes) avant de lever une exception. En pratique, une réponse valide est obtenue
en un ou deux essais ; le troisième n'a jamais été nécessaire sur les datasets testés.

\subsection{Persistance des checkpoints}

LangGraph utilise un \texttt{SqliteSaver} comme checkpointer. À chaque nœud exécuté,
l'état complet du pipeline est sérialisé. Si le processus est interrompu entre deux
validations humaines, la reprise se fait en réinstanciant le pipeline avec le même
\texttt{thread\_id}. L'état est restauré depuis le dernier checkpoint, sans nécessiter
de ré-exécuter les étapes précédentes.

\subsection{Versionnage des prompts}

Les prompts sont centralisés dans \texttt{config/prompt\_templates.yaml} avec un champ
\texttt{version}. Chaque rapport JSON inclut ce numéro de version, ce qui permet de
lier un résultat à la version exacte des instructions données au LLM. Une modification
de prompt est ainsi visible lors de l'analyse de variation de qualité entre deux runs
sur le même dataset.

% ============================================================
\section{Interfaces}
% ============================================================

Le pipeline est accessible selon trois modes. La CLI (\texttt{run.py}) expose l'ensemble
des paramètres via \texttt{argparse} et permet l'automatisation sur serveur. L'API REST
FastAPI expose trois endpoints pour l'intégration dans des systèmes externes ; chaque
session est identifiée par un UUID, et les checkpoints SQLite permettent de retrouver
l'état entre les requêtes HTTP. L'interface Streamlit (\texttt{app\_streamlit.py}) est
le mode le plus adapté à un usage interactif : elle reconstruit le flux human-in-the-loop
sous forme de formulaires, affiche les propositions de chaque agent avec leur niveau de
confiance et leur source, et permet de télécharger le rapport final en JSON.

% ============================================================
\section{Tests}
% ============================================================

\subsection{Stratégie}

Le projet contient 43 tests unitaires couvrant exclusivement les fonctions déterministes,
sans aucun appel LLM. Ce périmètre est délibéré. Les fonctions déterministes sont les
plus critiques : un bug dans la détection de bimodalité affecte silencieusement tous les
datasets qui passent par l'agent d'outliers. Elles sont aussi les seules qui peuvent être
testées de façon reproductible sans budget de tokens ni tolérance au non-déterminisme.

\begin{table}[H]
\centering
\begin{tabular}{lll}
\toprule
\textbf{Module} & \textbf{Fonction testée} & \textbf{Tests} \\
\midrule
\texttt{base\_agent} & \texttt{\_validate\_target} & 10 \\
\texttt{base\_agent} & \texttt{\_audit\_dtypes} & 7 \\
\texttt{base\_agent} & \texttt{\_detect\_trivial\_columns} & 7 \\
\texttt{outlier\_agent} & \texttt{\_bimodality\_coefficient} & 5 \\
\texttt{outlier\_agent} & \texttt{\_choose\_method} & 5 \\
\texttt{outlier\_agent} & \texttt{\_choose\_treatment} & 9 \\
\midrule
\multicolumn{2}{l}{\textbf{Total}} & \textbf{43} \\
\bottomrule
\end{tabular}
\caption{Couverture des tests unitaires.}
\end{table}

% ============================================================
\section{Benchmark sur 6 datasets}
% ============================================================

Le notebook \texttt{benchmark\_6\_datasets.ipynb} évalue le pipeline sur six datasets
publics issus de seaborn-data, choisis pour couvrir des scénarios contrastés.

\begin{table}[H]
\centering
\small
\begin{tabular}{llrlp{5cm}}
\toprule
\textbf{Dataset} & \textbf{Tâche} & \textbf{Lignes} & \textbf{Split} & \textbf{Principales difficultés} \\
\midrule
Titanic & classification & 891 & stratified & NaN (age, deck), colonnes triviales \\
Diamonds & régression & 53\,940 & auto & Catégoriques ordinales, fort volume \\
Penguins & classification & 344 & stratified & Multi-classe, NaN biologiques \\
MPG & régression & 398 & auto & Colonne texte proche d'un ID, NaN (horsepower) \\
Tips & régression & 244 & auto & Petit dataset, toutes colonnes catégoriques \\
Healthexp & régression & 540 & temporal & Série temporelle, split par année \\
\bottomrule
\end{tabular}
\caption{Datasets du benchmark.}
\end{table}

\begin{figure}[H]
\centering
\begin{tikzpicture}
\begin{axis}[
  ybar,
  bar width=18pt,
  width=\linewidth,
  height=6cm,
  symbolic x coords={Titanic,Diamonds,Penguins,MPG,Tips,Healthexp},
  xtick=data,
  ymin=0, ymax=100,
  ylabel={Score qualité /100},
  xlabel={Dataset},
  nodes near coords,
  nodes near coords align={vertical},
  grid=major,
  grid style={dashed, gray!30},
  every node near coord/.append style={font=\scriptsize},
  x tick label style={font=\small},
]
\addplot[fill=codeblue!60, draw=codeblue] coordinates {
  (Titanic,78) (Diamonds,85) (Penguins,80) (MPG,72) (Tips,82) (Healthexp,74)
};
\end{axis}
\end{tikzpicture}
\caption{Scores de qualité par dataset (valeurs issues d'une exécution représentative).}
\label{fig:scores}
\end{figure}

Les scores les plus bas, MPG et Healthexp, s'expliquent par des situations que le pipeline
gère moins bien. Dans MPG, la colonne \texttt{name} est du texte libre avec une cardinalité
proche du nombre de lignes --- elle ressemble à un identifiant sans en être un formellement,
ce qui complique sa classification automatique. Dans Healthexp, la structure temporelle
limite les options d'imputation croisée : les valeurs manquantes ne peuvent pas être
comblées par des observations futures sans introduire une fuite temporelle.

% ============================================================
\section{Limites et perspectives}
% ============================================================

\subsection{Limites identifiées}

\paragraph{Détection d'outliers univariée.}
L'IQR et le Z-score traitent chaque colonne indépendamment. Un point peut être dans les
bornes de chaque variable individuellement tout en étant un outlier multivariant ---
taille de 210~cm associée à un poids de 50~kg, par exemple. Des méthodes comme la
distance de Mahalanobis ou l'Isolation Forest comblent cette lacune, mais impliquent un
coût calculatoire plus élevé et des choix de paramétrisation supplémentaires.

\paragraph{Prompts en français.}
Les prompts sont rédigés en français. Pour un dataset dont les en-têtes de colonnes sont
en anglais, l'inférence de domaine peut être sous-optimale : le modèle doit jongler entre
la langue des instructions et la langue des données, ce qui dégrade parfois la qualité
des inférences sémantiques.

\paragraph{Support limité à CSV.}
Les formats Parquet, Excel, JSON et les connexions SQL nécessitent une extension du
nœud \texttt{analyze}. Il n'y a pas de difficulté technique particulière ; c'est une
question de périmètre d'implémentation.

\paragraph{Feature engineering informatif uniquement.}
Les propositions de features dérivées du LLM ne sont pas exécutées automatiquement.
Ce choix est conservateur mais justifié : exécuter du code Python proposé par un modèle
de langage sans validation formelle introduit un vecteur d'injection arbitraire.
Une exécution contrôlée (sandbox, validation statique du code) lèverait cette
contrainte.

\subsection{Perspectives}

Les extensions les plus directes sont le support de formats additionnels, l'exécution
contrôlée des features engineerées via sandbox, et le support multi-provider LLM
(Anthropic Claude, Mistral, modèles locaux via Ollama) --- les clés API sont déjà
présentes dans le fichier d'environnement en tant que variables commentées. La
calibration du score de qualité sur des benchmarks publics (OpenML, Kaggle) permettrait
de valider sa pertinence au-delà des six datasets actuels. Enfin, une détection
d'outliers multivariée (Isolation Forest, LOF) comblerait la principale limite
statistique du pipeline.

% ============================================================
\section{Conclusion}
% ============================================================

Le projet part d'un constat simple : la majorité des erreurs de preprocessing ne sont
pas dues à une incompréhension des concepts, mais à la difficulté de maintenir une
discipline rigoureuse sur des dizaines de micro-décisions répétitives. L'automatisation
résout ce problème pour les cas non ambigus ; le LLM le résout pour les cas contextuels ;
les interruptions humaines réservent la validation aux décisions à enjeu.

Ce qui distingue cette approche d'un prototype est la rigueur des garanties offertes :
étanchéité train/test structurellement impossible à violer, validation systématique des
sorties LLM avec boucle de réparation, persistance des checkpoints permettant la reprise
après interruption, et versionnage des prompts liant chaque résultat à la version exacte
des instructions utilisées. Ces propriétés rendent le pipeline reproductible et
auditable, deux conditions nécessaires pour l'utiliser dans un contexte professionnel
sur des datasets de structure variable.

\vspace{1cm}
\noindent\rule{\linewidth}{0.4pt}

\appendix

\section*{Annexe : Fichiers de configuration}

\subsection*{model\_config.yaml}

\begin{lstlisting}[language=Python, caption={Configuration du modèle LLM.}]
model: gpt-4o-mini
temperature: 0.1
max_tokens: 2000
max_retries: 3
retry_base_delay: 2.0
cache_enabled: true
cache_db_path: data/cache/llm_cache.db
\end{lstlisting}

\subsection*{prompt\_templates.yaml (extrait)}

\begin{lstlisting}[language=Python, caption={En-tête du fichier de prompts.}]
version: "1.1"

infer_domain: |
  Dataset: {dataset_name}
  Colonnes: {columns}
  Echantillon: {sample}
  Statistiques: {stats}
  ...

propose_transformations: |
  Colonnes ambiguës: {ambiguous_columns}
  Contexte domaine: {domain_context}
  ...
\end{lstlisting}

\end{document}
